\section{Lagrangian linearity}\label{sec:lagrangian}

By theorem~\ref{thm:slice}, each slice is a passive scalar carried by the velocity $\boldsymbol U_u=\uE+su\,\bhat$ of (\ref{eq:Uu}). The $\boldsymbol E\times\boldsymbol B$ part of this velocity is the same for all slices, and the streaming part is directed along the perturbed field. When the field lines are frozen into the $\boldsymbol E\times\boldsymbol B$ flow, the flow moves the field lines and the slices together, so that in the frame of the flow each slice streams along field lines that do not move. In this section we make this statement exact, list the hypotheses it needs, and derive its consequences for the quantities that are used to diagnose transfer in velocity space: the Hermite spectrum and flux, the slice correlations, and the distribution of the compressive fluctuations over wavenumbers along the field.

\subsection{Earlier work}\label{lag:earlier}

The linearity of the KRMHD kinetic equation in the Lagrangian frame of the Alfv\'enic flow is due to S09 (\S6.3). They transformed the kinetic equation (their eq.~157) to the Lagrangian frame of the flow $\uE$, used the Cauchy solution for the perturbed magnetic field to show that in this frame $\bhat\cdot\nabla$ is the derivative along the field lines taken at $t=0$ (their eqs~214--216), and concluded that the kinetic equation is linear in that frame. In their \S6.2.4 they remarked that the ballistic response $\exp(-\ii k_\parallel v_\parallel t)$ is also the nonlinear solution if $t$ and $k_\parallel$ are interpreted as Lagrangian variables. In \S6.3 they drew the consequence that the density and field-strength fluctuations do not cascade to small scales along the moving field lines, and that a cascade in $k_z$ appears only because the Alfv\'enic turbulence deforms the field lines in the perpendicular direction. \citet[\S\S2.6 and~4.1]{Kunz2015} restated the result for pressure-anisotropic plasmas: in the Lagrangian frame of the Alfv\'enic fluctuations, the compressive fluctuations obey a one-dimensional linear equation, which may be solved independently of the Alfv\'enic turbulence. For a passive scalar in MHD turbulence, \citet[\S2.3.1]{LithwickGoldreich2001} had noted that without dissipation the scalar would be frozen to the field lines and could not cascade along them, and argued that dissipation, which cannot be neglected, gives it a parallel cascade. The lemma beneath these statements is the classical one that $\boldsymbol B/\rho$ is Lie-dragged by the flow, which gives Cauchy's solution of the ideal induction equation \citep[example~3]{Webb2014}.

This section adds three things to these statements: the solution operator, written out in theorem~\ref{thm:frozen}; the hypotheses, stated explicitly together with the terms that violate them (section~\ref{lag:breaking}); and the use of the fact that the $\boldsymbol E\times\boldsymbol B$ flow preserves volume, which turns the theorem into exact statements about volume-integrated statistics (corollaries~\ref{cor:singlepoint}--\ref{cor:kpar}). None of the works cited above considers the Hermite spectrum or echoes.

\subsection{The frozen-line theorem}\label{lag:theorem}

We treat one compressive branch at a time and write its state at the point $\boldsymbol r$ as $g(\boldsymbol r,t)$. In the continuum, $g$ is the function $G(\boldsymbol r,\cdot,t)$ of $u$ and $L$ is the operator (\ref{eq:L}); in a truncated hierarchy, $g=(g_0,\dots,g_M)$ is the vector of Hermite coefficients and $L$ stands for the ladder matrix $A$. We write a collision or hypercollision term as $-\mathcal Cg$, as in (\ref{eq:sliceforcing}), so that a damping operator $\mathcal C$ is dissipative in the sense of (H2) below; GANDALF's hypercollisions, for example, are the diagonal matrix with entries $\nu(m/M)^n$ for $m\ge2$ and zero for $m=0,1$. The kinetic equation (\ref{eq:branch}) then reads
\begin{equation}\label{lag:kin}
  \partial_tg+\{\Phi,g\}+s\,\gradpar(Lg)=-\mathcal Cg .
\end{equation}
The Laguerre sectors with $l\ge2$ obey the same equation with $\Lambda=\infty$. All branches and sectors are advected by the same Alfv\'enic fields, so the results of this section hold for each of them with the same relabelling and apply to the whole compressive distribution at once; for the slices this is (\ref{eq:sameflow}).

We write $D_t=\partial_t+\{\Phi,\cdot\}$ for the convective derivative, $\gradpar(t)=\partial_z+\{\Psi(\cdot,t),\cdot\}$ for the derivative along the field at time $t$, $\langle\cdot\rangle$ for the average over the periodic box $\mathbb T^3$, and $\boldsymbol X_t$ for the flow map of $\uE$, defined by $\partial_t\boldsymbol X_t(\boldsymbol\xi)=\uE(\boldsymbol X_t(\boldsymbol\xi),t)$ and $\boldsymbol X_0(\boldsymbol\xi)=\boldsymbol\xi$. The point $\boldsymbol\xi$ labels the fluid element that is at $\boldsymbol X_t(\boldsymbol\xi)$ at time $t$. The relabelling operator
\begin{equation*}
  \mathcal X_th=h\circ\boldsymbol X_t^{-1}
\end{equation*}
takes a function $h$ of the labels to the function of position whose value at $\boldsymbol X_t(\boldsymbol\xi)$ is $h(\boldsymbol\xi)$. The residual of the flux equation,
\begin{equation*}
  R=\partial_t\Psi+\{\Phi,\Psi\}-\partial_z\Phi ,
\end{equation*}
is proportional to the parallel electric field of the Alfv\'enic fluctuations, and the first of equations (\ref{eq:rmhd}) is the statement $R=0$. The hypotheses are
\begin{equation}\label{eq:hyp}
\begin{minipage}{0.86\linewidth}
\begin{list}{}{\setlength{\leftmargin}{2.9em}\setlength{\labelwidth}{2.5em}\setlength{\labelsep}{0.4em}\setlength{\itemsep}{0.6ex}\setlength{\parsep}{0pt}\setlength{\topsep}{0pt}}
\item[(H1)] \emph{Ideal induction.} On $[0,T]$ the residual $R$ depends on $z$ and $t$ only. Nothing is assumed about the vorticity equation: $\Phi$ may be forced, viscous or prescribed.
\item[(H2)] \emph{Kinetic equation.} The branch has $\Lambda\notin[0,1]$, and $g$ obeys (\ref{lag:kin}), where $\mathcal C$ is a linear operator on the velocity variable that acts at each point $\boldsymbol r$ separately, with a kernel that depends neither on $\boldsymbol r$ nor on the fields (it may depend on $t$), and $\mathcal C$ is either bounded or dissipative, where dissipative means that $\mathrm{Re}\int\fe{\bar f}{\mathcal Cf}\,\dd^3r\ge0$ for every $f$ in its domain. There is no spatial dissipation of $g$ and no source.
\item[(H3)] \emph{Smoothness.} $\Phi$, $\Psi$ and all their spatial derivatives are continuous functions of $(\boldsymbol r,t)$ on $\mathbb T^3\times[0,T]$, except that $\Phi$ may jump at finitely many times (impulsive kicks); that is, $[0,T]$ lies within the interval on which the Alfv\'enic solution exists and is smooth.
\end{list}
\end{minipage}
\end{equation}
A residual that depends on $z$ and $t$ alone can be removed by adding a function of $z$ and $t$ to $\Psi$, which changes neither $\boldsymbol b_\perp$ nor $\gradpar$, so (H1) is in effect the flux equation of (\ref{eq:rmhd}), with any forcing or dissipation acting on $\Phi$. Under (H3), $\uE$ is divergence-free, smooth in $\boldsymbol r$ and continuous in $t$, so $\boldsymbol X_t$ is a volume-preserving diffeomorphism of $\mathbb T^3$ that is continuously differentiable in $t$. Continuity in $t$ suffices, so white-noise forcing of $\Phi$, which gives fields that are continuous but not differentiable in time, is allowed. Kicks of $\Phi$ at finitely many times are also allowed: $\Psi$ and $\boldsymbol X_t$ are continuous across a kick, so the theorem applies on each interval between kicks and the relabellings compose. We write $\gradpar(0)$ for the derivative along the initial field and define the frozen generator
\begin{equation*}
  L_0=s\,\gradpar(0)\,L+\mathcal C .
\end{equation*}
Its solution operator $e^{-tL_0}$, or the propagator if $\mathcal C$ depends on $t$, solves the linear problem in which the field lines keep their initial shape and there is no flow.

\begin{theorem}[Frozen-line theorem]\label{thm:frozen}
Assume (\ref{eq:hyp}), and let $g$ be a classical solution of (\ref{lag:kin}) on $[0,T]$, with $g$ and $Lg$ continuously differentiable in $(\boldsymbol r,t)$. Then
\begin{equation}\label{eq:frozen}
  g(t)=\mathcal X_t\,e^{-tL_0}\,g(0),\qquad L_0=s\,\gradpar(0)\,L+\mathcal C,\qquad 0\le t\le T .
\end{equation}
Equivalently, the function of the labels $\hat g(\boldsymbol\xi,t)=g(\boldsymbol X_t(\boldsymbol\xi),t)$ obeys $\partial_t\hat g+s\,\gradpar(0)(L\hat g)=-\mathcal C\hat g$: in the labels, $g$ evolves as in the static field $\bhat(0)$, with no flow.
\end{theorem}

\begin{proof}
With $\bhat=\zhat+\zhat\times\nabla_\perp\Psi$ one has $\uE\times\bhat=\nabla_\perp\Phi+\{\Phi,\Psi\}\,\zhat$, and therefore
\begin{equation*}
  \partial_t\bhat-\nabla\times(\uE\times\bhat)=\zhat\times\nabla R .
\end{equation*}
The right-hand side vanishes if and only if $R$ depends on $z$ and $t$ only, so (H1) is equivalent to the ideal induction equation for the full field $\bhat$. Because $\nabla\cdot\uE=\nabla\cdot\bhat=0$, the flow map preserves volume, and the ideal induction equation has Cauchy's solution
\begin{equation*}
  \bhat\big(\boldsymbol X_t(\boldsymbol\xi),t\big)=\frac{\partial\boldsymbol X_t}{\partial\boldsymbol\xi}\,\bhat(\boldsymbol\xi,0) .
\end{equation*}
For a differentiable function $h$ of the labels, the chain rule then gives $\bhat(t)\cdot\nabla(h\circ\boldsymbol X_t^{-1})=[\bhat(0)\cdot\nabla h]\circ\boldsymbol X_t^{-1}$, that is,
\begin{equation}\label{lag:intertwine}
  \gradpar(t)\,\mathcal X_t=\mathcal X_t\,\gradpar(0) .
\end{equation}
In RMHD variables the same relation follows from the identity
\begin{equation}\label{lag:commutator}
  [D_t,\gradpar]\,f=\{R,f\},
\end{equation}
which holds for every $f$ and follows from the Jacobi identity for the bracket. Under (H1) its right-hand side vanishes. For a time-independent $h$ we have $D_t(\mathcal X_th)=0$, hence $D_t(\gradpar(t)\mathcal X_th)=0$, so $\gradpar(t)\mathcal X_th$ is carried by the flow from its initial value $\gradpar(0)h$, which is (\ref{lag:intertwine}).

Now write $g(t)=\mathcal X_t\hat g(t)$. By the chain rule and (\ref{lag:kin}), $\partial_t\hat g=(D_tg)\circ\boldsymbol X_t=-[s\gradpar(t)(Lg)+\mathcal Cg]\circ\boldsymbol X_t$. Because $L$ and $\mathcal C$ act on the velocity variable at each point, with kernels that do not depend on $\boldsymbol r$, they commute with composition by $\boldsymbol X_t$, and by (\ref{lag:intertwine}) $[\gradpar(t)(Lg)]\circ\boldsymbol X_t=\gradpar(0)(L\hat g)$. Hence $\hat g$ solves the frozen problem $\partial_t\hat g=-L_0\hat g$ with $\hat g(0)=g(0)$, whose coefficients do not depend on time apart from any time dependence of $\mathcal C$. The solution of this problem is unique (see below), so $\hat g(t)=e^{-tL_0}g(0)$, which is (\ref{eq:frozen}).
\end{proof}

The proof uses the induction equation, the fact that $L$ and $\mathcal C$ act at each point with coefficients that do not depend on position, and nothing else. The branch operator has this property because the field relations of KRMHD are algebraic (section~\ref{sec:model}). The vorticity equation never enters. For any Alfv\'enic fields that satisfy (H1) and (H3), the problem in the labels is linear, with coefficients that do not depend on time unless $\mathcal C$ does; the flow only relabels the points.

Without collisions, (\ref{eq:frozen}) decouples into slices. The map $\mathcal A$ to the unit-weight amplitudes (\ref{eq:aunit}) acts on $u$ alone, with a kernel that does not depend on position, so it commutes with $\mathcal X_t$ and with $\gradpar(0)$, and it takes $L$ to multiplication by $u$ (theorem~\ref{thm:spectrum}). The amplitudes therefore obey
\begin{equation}\label{lag:sliceform}
  a_u(t)=\mathcal X_t\,T_{sut}\,a_u(0),\qquad T_\ell=e^{-\ell\gradpar(0)} .
\end{equation}
Because $\bhat(0)$ is smooth and divergence-free and $\bhat(0)\cdot\zhat=1$, the operator $T_\ell$ moves a function a distance $\Delta z=\ell$ along the field lines of $\bhat(0)$, and the $T_\ell$ form a group of unitary operators on $L^2(\mathbb T^3)$, with $T_\ell T_{\ell'}=T_{\ell+\ell'}$. Each slice is carried rigidly along the initial field lines at the speed $su$, and the flow then relabels the points. This makes exact, for the slices, the remark of S09 (\S6.2.4) that the ballistic response is also the nonlinear solution in Lagrangian variables. In a truncated hierarchy, (\ref{lag:sliceform}) holds for the $M+1$ slices of theorem~\ref{thm:trunc}, with $u$ replaced by the eigenvalues $u_j$ of $A$. With collisions, $\mathcal C$ and $L$ do not commute in general, the exponential in (\ref{eq:frozen}) does not split into slices, and the slices remain coupled through $\mathcal C$ as in (\ref{eq:sliceforcing}); the theorem still removes the flow.

The algebra of the proof does not involve $\Lambda$. The condition $\Lambda\notin[0,1]$ in (H2) makes the frozen problem well posed. For $\Lambda\notin[0,1]$ the map $\mathcal A:G\mapsto a$ of (\ref{eq:aunit}) is unitary from $(\mathcal W,\fe{\cdot}{\cdot})$ onto $L^2(\mathbb R,\dd u)$ (theorem~\ref{thm:spectrum}), so $\int\fe{\bar f}{f}\,\dd^3r$ is the square of a norm. The streaming part $s\gradpar(0)L$ of $L_0$ is antisymmetric for the product $\int\fe{\bar\cdot}{\cdot}\,\dd^3r$, because $L$ is symmetric in the free-energy product and $\gradpar(0)$, a derivative along a divergence-free field, is antisymmetric in $L^2(\mathbb T^3)$. The collisionless $e^{-tL_0}$ therefore preserves the free energy $\int\fe{g}{g}\,\dd^3r$, and a dissipative $\mathcal C$, such as hypercollisions or the Lenard--Bernstein operator, can only decrease it. For such a $\mathcal C$, or for a bounded one with Gronwall's inequality, the same estimate applied to the difference of two solutions shows that the frozen problem has at most one solution for given initial data. For $0<\Lambda<1$, $L$ has the eigenvalues $\pm\ii\gamma$ with $\gamma>0$ (section~\ref{sec:slices}). In a uniform field, a Fourier mode with wavenumber $k$ along the field whose velocity profile is the corresponding eigenfunction then grows at the rate $s|k|\gamma$, which has no upper bound, so the frozen problem is ill posed and (\ref{eq:frozen}) holds only for solutions that exist. The KRMHD branches have $\Lambda^+>1$ and $\Lambda^-<0$ for all $\beta_i$ and $\tau/Z$ (section~\ref{sec:model}), and the Laguerre sectors have $\Lambda=\infty$, so the condition holds throughout KRMHD.

Hypothesis (H3) restricts the scope of the theorem. Whether solutions of the ideal three-dimensional RMHD equations stay smooth for all time is not known, so the theorem is a statement about the interval on which they do. A numerical test of the theorem is limited further, to times at which the Alfv\'enic fields themselves are resolved. In RK4 integrations of GANDALF's right-hand side \citep{Kanekar2025} with $\mathrm{rms}\,|\nabla_\perp\Phi|=\mathrm{rms}\,|\nabla_\perp\Psi|=1.5$, Cauchy's formula, which involves only the Alfv\'enic fields, is satisfied at $t=0.5$ only to a relative error of 0.25, 0.27 and 0.17 on $32^3$, $48^3$ and $64^3$ grids (the largest difference between the two sides at 40 Lagrangian markers, relative to the largest component of $\boldsymbol b_\perp$ there).

\subsection{Single-point statistics}\label{lag:singlepointsec}

Let $O$ be a function of the local state of $g$ and of finitely many of its derivatives along the field. The local state is all the velocity-space information at one point: the function $G(\boldsymbol r,\cdot)$, or the coefficients $g_m(\boldsymbol r)$, of every branch and Laguerre sector. For a field $f$ on $\mathbb T^3$ with values in velocity space we write
\begin{equation}\label{lag:Odef}
  \mathcal O_t[f]=\int O\big(f,\gradpar(t)f,\gradpar(t)^2f,\dots\big)\,\dd^3r ,
\end{equation}
where the subscript records the time at which the derivatives along the field are taken, and we call $\mathcal O_t[g(t)]$ a single-point statistic. Since $\boldsymbol X_t$ preserves volume, $\int\mathcal X_th\,\dd^3r=\int h\,\dd^3\xi$ for every integrable $h$, and by (\ref{lag:intertwine}) $\gradpar(t)^k\mathcal X_t=\mathcal X_t\gradpar(0)^k$ for every $k$. Together these give
\begin{equation}\label{lag:Oinv}
  \mathcal O_t[\mathcal X_ty]=\mathcal O_0[y]
\end{equation}
for every $y$ for which either side is defined.

\begin{corollary}[Single-point statistics]\label{cor:singlepoint}
Assume (\ref{eq:hyp}), and let $h(t)=e^{-tL_0}g(0)$ be the solution of the frozen problem. Then, for every single-point statistic (\ref{lag:Odef}) and $0\le t\le T$,
\begin{equation}\label{lag:singlepoint}
  \mathcal O_t[g(t)]=\mathcal O_0[h(t)] .
\end{equation}
Every volume-integrated single-point statistic of $g$ therefore takes the value it has in the linear problem along the field lines frozen at $t=0$, with no flow.
\end{corollary}

\begin{proof}
By theorem~\ref{thm:frozen}, $g(t)=\mathcal X_th(t)$; apply (\ref{lag:Oinv}) with $y=h(t)$.
\end{proof}

The corollary covers the Hermite spectrum $W(m,t)=p_m\langle g_m^2\rangle$ and the free energy; the variances of $\delta n$ and $\delta B_\parallel$; the slice energies and the full slice correlation matrix $C(u,u';t)=\langle a_ua_{u'}\rangle$, because the slice amplitudes are linear functions of the local state; the collisional dissipation rate $\varepsilon_\nu$, a volume average of $\fe{g}{\mathcal Cg}$; the Hermite flux $\Pi_{m+1/2}$ written with the full $\gradpar$, as in (\ref{eq:fluxCD}); and the one-point probability distributions of the $g_m$, with all their moments and quantiles. It does not cover spectra in $k_\perp$ or $k_z$, or fluxes resolved in $k_\perp$ shells, which involve Fourier filters that do not commute with $\mathcal X_t$; the flux computed with $\partial_z$ in place of $\gradpar$, which involves a derivative that is not taken along the field; or two-point statistics at a fixed Eulerian separation.

Without collisions, (\ref{lag:sliceform}), volume preservation and the unitarity of $T_\ell$ give the slice correlations in closed form,
\begin{equation}\label{eq:frozencorr}
  C(u,u';t)=\big\langle a_u(0)\,T_{s(u'-u)t}\,a_{u'}(0)\big\rangle .
\end{equation}
The correlation of two slices at time $t$ is the correlation of the initial slices at the field-aligned separation $s(u'-u)t$ along the initial field lines. Under (\ref{eq:hyp}), (\ref{eq:frozencorr}) integrates the two-slice law of proposition~\ref{prop:twoslice} in closed form, and phase mixing is the decay of the initial correlations with separation along the frozen field lines. The Hermite spectrum is a quadratic form in the slice correlations (section~\ref{sec:decorrelation}), so it is fixed by (\ref{eq:frozencorr}). In a truncated hierarchy, where $g_m=\sum_jR_{mj}a_j$ with $R_{mj}=r_m(u_j)/\sqrt{n_j}$ and $n_j=\sum_mp_mr_m(u_j)^2$ (section~\ref{sec:decorrelation}),
\begin{equation*}
  W(m,t)=p_m\sum_{j,l}R_{mj}R_{ml}\,\big\langle a_j(0)\,T_{s(u_l-u_j)t}\,a_l(0)\big\rangle .
\end{equation*}

We tested corollary~\ref{cor:singlepoint} with RK4 integrations of GANDALF's right-hand side at $M=8$, $\Lambda=\sqrt5$ and $s=1$, from random fields with $\mathrm{rms}\,|\nabla_\perp\Phi|=\mathrm{rms}\,|\nabla_\perp\Psi|=1$, for which the $\boldsymbol E\times\boldsymbol B$ term in the equation for $g_0$ is 2.9 times as large, in r.m.s.\ value, as the streaming term $s\,\partial_z(Ag)_0$. The initial fields are the same physical fields at every resolution, so a statistic covered by the corollary must converge to its value for the frozen-line prediction $h(t)$ as the grid is refined, and one that is not covered need not. For the Hermite spectrum we use the relative difference $\sum_m|\langle g_m^2\rangle-\langle g_m^2\rangle_{\mathrm{fr}}|/\sum_m\langle g_m^2\rangle_{\mathrm{fr}}$, where the subscript denotes the frozen-line prediction. At $t=0.062$ it is $3.4\times10^{-4}$, $2.0\times10^{-5}$, $8.9\times10^{-6}$ and $2.2\times10^{-7}$ on $16^3$, $24^3$, $32^3$ and $48^3$ grids. The slice correlation matrix, the field-aligned flux and the fourth moments $\langle g_m^4\rangle$ converge in the same way, whereas the spectrum $W(k_\perp)$ differs from its frozen-line counterpart by 0.21, and the $k_z$-based flux by 0.069, at every resolution. The same holds in runs with GANDALF's hypercollisions ($\nu=20$, $n=4$), in which the collisional dissipation rate $\varepsilon_\nu$ also converges. The linear prediction with $\partial_z$ in place of $\gradpar(0)$, which has neither the flow nor the field-line perturbation, misses the Hermite spectrum by 6\% and the flux by 44\%. In a pointwise check along 40 Lagrangian markers, with different random fields of the same amplitude, the largest difference between $g(\boldsymbol X_t(\boldsymbol\xi),t)$ and $h(\boldsymbol\xi,t)$ at $t=0.1$ is $3.8\times10^{-3}$, $3.4\times10^{-4}$, $5.8\times10^{-5}$ and $1.6\times10^{-6}$ of $\max|h|$ on the same four grids, against 0.18 for the comparison without relabelling. At $t=0.25$, about three eddy times, the differences in the Hermite spectrum, the slice correlations, the flux and the quantiles still decrease from $24^3$ to $48^3$, while those in the fourth moments, and in $\varepsilon_\nu$ in the runs with hypercollisions, are not monotone in the resolution; by then the passive field has reached the grid scale. Through GANDALF's own Lawson-RK4 time step on a $32^3$ grid, the difference in the Hermite spectrum at $t=0.064$ is $9.7\times10^{-6}$, $9.3\times10^{-6}$ and $9.9\times10^{-6}$ for $\Delta t=0.002$, $0.001$ and $0.0005$, which is the spatial error of that grid; the IMEX scheme is discussed in appendix~\ref{app:timestepping}. Section~\ref{sec:numerics} and figure~\ref{fig:frozen} give the full comparison.

\subsection{Forcing}\label{lag:forcingsec}

With a source $F(\boldsymbol r,t)$ added to the right-hand side of (\ref{lag:kin}), the same change of variables gives $\partial_t\hat g=-L_0\hat g+\mathcal X_t^{-1}F(t)$, and Duhamel's formula in the labels reads
\begin{equation}\label{lag:duhamel}
  g(t)=\mathcal X_t\Big[e^{-tL_0}g(0)+\int_0^te^{-(t-t')L_0}\,\mathcal X_{t'}^{-1}F(t')\,\dd t'\Big] .
\end{equation}
By (\ref{lag:intertwine}), $\mathcal X_{t'}e^{-(t-t')L_0}\mathcal X_{t'}^{-1}=e^{-(t-t')L_{t'}}$ with $L_{t'}=s\,\gradpar(t')\,L+\mathcal C$, so the source injected at $t'$ contributes $\mathcal X_t\mathcal X_{t'}^{-1}e^{-(t-t')L_{t'}}F(t')$ to $g(t)$. Each injection phase-mixes along the field lines frozen at its injection time and is then relabelled by the flow between $t'$ and $t$. (For a time-dependent $\mathcal C$ the exponentials become propagators.)

For a single realisation, a quadratic statistic of $g(t)$ contains cross terms between injections at different times, and these involve the relabelling between the injections. For forcing that is white in time they vanish in the mean. To state this precisely, let $(\mathfrak F_t)_{t\ge0}$ be a filtration, the information available at time $t$, to which $g(0)$, the Alfv\'enic fields and the forcing are adapted. Let the source be $F\,\dd t=\dd Y$, where $Y$ is a Wiener process with spatially smooth realisations and covariance operator $\Sigma$ in position and velocity, whose increment $Y(t'')-Y(t')$ is independent of $\mathfrak F_{t'}$ for $t''>t'$. We call such forcing non-anticipating. The Alfv\'enic fields may be driven by a noise that is correlated with $Y$ at the same instant: the flow responds to such a kick only after it, while the integrand of (\ref{lag:duhamel}) at $t'$ depends only on the flow before $t'$. Forcing for which the flow at time $t'$ depends on the forcing of $g$ after $t'$ is excluded.

\begin{corollary}[White forcing]\label{cor:forcing}
Let the forcing be non-anticipating, let (H1) and (H3) of (\ref{eq:hyp}) hold for every realisation, and let $g$ obey (\ref{lag:kin}) with the source $F$ and otherwise satisfy (H2). Let $\mathcal O_t$ be a single-point statistic (\ref{lag:Odef}) whose integrand is a quadratic form. Then
\begin{equation}\label{lag:meanforced}
  \mathrm E\,\mathcal O_t[g(t)]=\mathrm E\,\mathcal O_0\big[e^{-tL_0}g(0)\big]+\int_0^t\mathrm E\,\mathcal O_{t'}\big[e^{-(t-t')L_{t'}}\chi\big]\,\dd t' ,
\end{equation}
where $\chi$ is a random field with zero mean and covariance $\Sigma$, independent of all other variables. In particular, the mean Hermite spectrum $\mathrm E\,W(m,t)$ is the linear response to the initial condition plus the sum of the linear responses to the individual injections, each phase-mixing along the field lines frozen at its own injection time, with no further effect of the flow.
\end{corollary}

\begin{proof}
In the labels the forced equation is the linear equation $\dd\hat g=-L_0\hat g\,\dd t+\mathcal X_t^{-1}\dd Y$. Its solution is the bracket in (\ref{lag:duhamel}), with the integral read as an It\^o integral, which is well defined because its integrand $e^{-(t-t')L_0}\mathcal X_{t'}^{-1}$ is fixed by $\mathfrak F_{t'}$. The noise does not depend on $g$ and $\boldsymbol X_t$ is differentiable in $t$, so the It\^o and Stratonovich readings of the forced equation coincide. By (\ref{lag:Oinv}), $\mathcal O_t[g(t)]=\mathcal O_0[\hat g(t)]$. The cross term between $e^{-tL_0}g(0)$ and the stochastic integral has zero mean, because the first is fixed by $\mathfrak F_0$ and the integral has zero mean given $\mathfrak F_0$. The It\^o isometry turns the mean of the remaining term into $\int_0^t\mathrm E\,\mathcal O_0[e^{-(t-t')L_0}\mathcal X_{t'}^{-1}\chi]\,\dd t'$. Finally, (\ref{lag:Oinv}) at time $t'$ gives $\mathcal O_0[\mathcal X_{t'}^{-1}y]=\mathcal O_{t'}[y]$ for every $y$, and with $y=\mathcal X_{t'}e^{-(t-t')L_0}\mathcal X_{t'}^{-1}\chi=e^{-(t-t')L_{t'}}\chi$ this is (\ref{lag:meanforced}).
\end{proof}

The same argument applies, with no limit to be taken, to forcing applied as kicks $g\to g+\Delta Y_n$ at times $t_n$, where each $\Delta Y_n$ has zero mean and is independent of everything that happened before $t_n$; an Alfv\'enic kick at the same instant is again allowed. The integral in (\ref{lag:meanforced}) is then replaced by the sum of $\mathrm E\,\mathcal O_{t_n}[e^{-(t-t_n)L_{t_n}}\Delta Y_n]$ over the kicks with $t_n\le t$. The forcing routines of GANDALF act in this way: once per time step they add an independent increment whose amplitude is proportional to $\Delta t^{-1/2}$.

The corollary concerns the mean: for a single realisation, $W(m,t)$ contains cross terms between injections at different times. It needs forcing that is white in time: for forcing correlated over a time $\tau_c$, the mean cross terms between injections less than about $\tau_c$ apart survive, and they contain the relabelling $\mathcal X_{t''}\mathcal X_{t'}^{-1}$ between the two injection times, that is, the flow. It holds for quadratic statistics only: a quartic statistic such as $\langle g_m^4\rangle$ contains products of the contributions of two injection times at the same point, which again involve that relabelling. Finally, each injection needs the hypotheses on $[t',t]$. Because (H1) excludes resistivity, the magnetic field is not dissipated, and the corollary concerns finite intervals on which the solution stays smooth, not the statistically steady state of a resistive simulation.

We checked corollary~\ref{cor:forcing} with RK4 integrations of GANDALF's right-hand side at $M=8$, $\Lambda=\sqrt5$ and $s=1$, in turbulence forced through $\Phi$ alone, with the same white-noise increment added to $z^+$ and $z^-$ at every step. Three impulsive injections into $g$ were made at $t=0$, $0.04$ and $0.08$; the eddy time is about $0.1$. At $t=0.16$ the relative difference between the Hermite spectrum of the response to each injection and the frozen-line prediction from its own injection time is $4.2\times10^{-4}$, $7.2\times10^{-4}$ and $2.3\times10^{-4}$ on a $48^3$ grid, down from $2.4\times10^{-3}$, $3.1\times10^{-3}$ and $1.0\times10^{-3}$ on $24^3$. For each of the $2^3$ sign patterns of the injections the cross terms are 5--11\% of $W$; the average over the patterns removes them identically. When the same pattern is injected twice with the same sign, which models forcing correlated over the gap between the injections, the cross term is 0.90, 0.69 and 0.49 times the diagonal part for gaps of 0.04, 0.08 and 0.12, and an estimate that omits the relabelling between the injections misses it by 3\%, 12\% and 38\%. In these runs the injections do not affect the flow. A sharper test applies kicks to $\Phi$ at the two injection times $t=0$ and $0.06$, each with the same r.m.s.\ gradient as the initial flow and with the random sign of the injection into $g$ made at that time, so that the flow after each injection depends on its sign. The relative difference $\sum_m|\Delta W(m)|/\sum_mW(m)$ between the sign-averaged free-energy spectrum $W(m)=p_m\langle g_m^2\rangle$ at $t=0.15$ and the averaged sum of the frozen-line responses is $3.4\times10^{-5}$, $2.3\times10^{-7}$, $1.6\times10^{-8}$ and $9.8\times10^{-11}$ on $16^3$, $24^3$, $32^3$ and $48^3$ grids, although the Eulerian response to the first injection changes by 26\% when the sign of the second kick is reversed. When the kick that carries the sign of the second injection is applied at $t=0$, before that injection, so that the forcing anticipates, the difference stays at $2.2\times10^{-4}$ at every resolution.

\subsection{Field-aligned wavenumbers}\label{lag:kparsec}

The theorem also fixes how the compressive fluctuations are distributed over wavenumbers along the field. For each $t$ the field $\bhat(t)$ is smooth and divergence-free, so the maps that move points along its field lines preserve volume, and the corresponding translations of functions form a strongly continuous group of unitary operators on $L^2(\mathbb T^3)$, with generator $-\gradpar(t)$. By Stone's theorem \citep[e.g.][]{Kato1966},
\begin{equation*}
  K_\parallel(t)=-\ii\,\gradpar(t)
\end{equation*}
is self-adjoint. We call its spectral variable $\kappa$ the field-aligned wavenumber and write $E_t$ for its spectral measure, so that $E_t(I)f$ is the part of $f$ with field-aligned wavenumbers in the Borel set $I\subset\mathbb R$; $E_t(I)$ acts on position, in the same way at every velocity. In a uniform field, $K_\parallel(t)=-\ii\partial_z$, $\kappa=k_z$, and $E_t(I)$ is the Fourier filter onto $k_z\in I$. When the field lines are tangled, $K_\parallel(t)$ need not have eigenfunctions, and the spectral measure takes the place of the Fourier transform along the field. With the norm $\|f\|^2=\int\fe{\bar f}{f}\,\dd^3r$, which is proportional to the free energy for real $f$, we define the field-aligned spectrum of the free energy,
\begin{equation*}
  \mu_t(I)=\big\|E_t(I)\,g(t)\big\|^2 ,
\end{equation*}
and the field-aligned spectrum of a slice in the same way with the norm of $L^2(\mathbb T^3)$.

\begin{corollary}[Field-aligned wavenumbers]\label{cor:kpar}
Assume (\ref{eq:hyp}). For each $t\in[0,T]$ the operator $K_\parallel(t)$ is self-adjoint, and for every Borel set $I\subset\mathbb R$
\begin{equation}\label{lag:kparsub}
  E_t(I)\,g(t)=\mathcal X_t\,e^{-tL_0}\,E_0(I)\,g(0) .
\end{equation}
The part of $g$ with field-aligned wavenumbers in $I$ therefore evolves, in the labels, from the part of $g(0)$ with field-aligned wavenumbers in $I$ alone: the spectral subspaces of $K_\parallel(0)$ are invariant under the dynamics in the labels, and the spectral measure of $g(t)$ with respect to $K_\parallel(t)$ equals that of $e^{-tL_0}g(0)$ with respect to $K_\parallel(0)$. Without collisions, $\mu_t(I)=\mu_0(I)$ for every $I$, and the field-aligned spectrum of each slice is constant in time. If $\mathcal C$ is dissipative, $\mu_t(I)$ does not increase.
\end{corollary}

\begin{proof}
Self-adjointness was shown above. By (\ref{lag:intertwine}) and the unitarity of $\mathcal X_t$ on $L^2(\mathbb T^3)$, $K_\parallel(t)=\mathcal X_tK_\parallel(0)\mathcal X_t^{-1}$, so $E_t(I)=\mathcal X_tE_0(I)\mathcal X_t^{-1}$. The projection $E_0(I)$ acts on position only and commutes with $\gradpar(0)$, so it commutes with $L_0$ and with $e^{-tL_0}$, and (\ref{eq:frozen}) gives (\ref{lag:kparsub}). Since $\mathcal X_t$ preserves the norm, $\mu_t(I)=\|E_0(I)\,e^{-tL_0}g(0)\|^2$, which is the spectral measure of $e^{-tL_0}g(0)$ with respect to $K_\parallel(0)$. Without collisions, $e^{-tL_0}$ acts on the slice at $u$ as $T_{sut}=e^{-\ii sutK_\parallel(0)}$, which multiplies the component at field-aligned wavenumber $\kappa$ by the phase $e^{-\ii su\kappa t}$; this gives the statements for each slice and, by theorem~\ref{thm:spectrum}, $\mu_t(I)=\mu_0(I)$. With collisions, let $y(t)=e^{-tL_0}E_0(I)g(0)$. The streaming part $s\gradpar(0)L$ of $L_0$ is antisymmetric for the product $\int\fe{\bar\cdot}{\cdot}\,\dd^3r$ (section~\ref{lag:theorem}), so $\dd\|y\|^2/\dd t=-2\,\mathrm{Re}\int\fe{\bar y}{\mathcal Cy}\,\dd^3r$, which is not positive for a dissipative $\mathcal C$.
\end{proof}

For a real $g$ the spectral measure is symmetric, $\mu_t(-I)=\mu_t(I)$, because $\gradpar(t)$ is real, so that $K_\parallel(t)$ anticommutes with complex conjugation; the parts $E_t(0,\infty)g$ and $E_t(-\infty,0)g$ are complex conjugates of each other. The statement that no part of $g$ changes the sign of its field-aligned wavenumber is therefore a statement about these two complex components: by (\ref{lag:kparsub}) they evolve separately, and neither leaks into the other. In a uniform field they are the parts of $g$ with $k_z>0$ and with $k_z<0$. Corollary~\ref{cor:kpar} gives a precise form of the conclusion of S09 (\S6.3) that the compressive fluctuations do not cascade to small scales along the moving field lines: under (\ref{eq:hyp}) and without collisions, the distribution of the free energy over field-aligned wavenumbers does not change in time.

The frozen-in shear shows the difference between the field-aligned and the Eulerian wavenumber. Take $\Phi=-Szy$, so that $\uE=Sz\,\hat{\boldsymbol x}$, and $\Psi=0$ at $t=0$. The flux equation is solved by $\Psi=-Sty$, so that $\boldsymbol b_\perp=St\,\hat{\boldsymbol x}$ and $\gradpar(t)=\partial_z+St\,\partial_x$; both perpendicular Laplacians vanish, and the pair is an exact solution of (\ref{eq:rmhd}), although not a periodic one. The flow map is $\boldsymbol X_t(\boldsymbol\xi)=(\xi_x+S\xi_zt,\,\xi_y,\,\xi_z)$, so a Fourier mode of the labels, $\exp[\ii(k_x\xi_x+k_y\xi_y+k_z\xi_z)]$, becomes at time $t$ a Fourier mode of position with
\begin{equation}\label{lag:shear}
  k_z(t)=k_z-Stk_x,\qquad k_\parallel(t)=k_z(t)+Stk_x=k_z ,
\end{equation}
where $k_\parallel(t)$ is its eigenvalue under $K_\parallel(t)$. The Eulerian wavenumber $k_z(t)$ changes sign at $t=k_z/(Sk_x)$, while the field-aligned wavenumber stays fixed, and in the labels the velocity-space evolution is that of a uniform field with wavenumber $k_z$ along it. A diagnostic that sorts the free energy by the sign of $k_z$ counts this mode as having reversed its parallel wavenumber, although nothing has changed along the field. Section~\ref{sec:decorrelation} works out the consequences for the decompositions of the Hermite flux that are based on $k_z$.

\subsection{What breaks the theorem}\label{lag:breaking}

Apart from the smoothness (H3), the proof uses (H1) and (H2) and nothing else about the dynamics, and without either of them the conclusion fails in general. A non-ideal term in the induction equation in general makes $R$ depend on $x$ and $y$, so that the commutator (\ref{lag:commutator}) does not vanish and the field lines are not carried by the flow. Resistivity or hyper-resistivity acting on $\Psi$, parallel dissipation applied to $z^\pm$, and any forcing of $\Psi$, for example independent forcing of $z^+$ and $z^-$, are terms of this kind. An operator on $g$ that does not commute with the relabelling breaks the theorem as well; dissipation of $g$ that depends on $k_\perp$ or on $k_z$, as in GANDALF, is in general such an operator. Viscosity or forcing acting on $\Phi$ alone does not break it, because the vorticity equation never enters the proof. The Fourier truncation of a spectral code is itself an operator that does not commute with the relabelling. The error it causes vanishes as the resolution increases at fixed $t$, but not uniformly in $t$, and it limits the time over which a given resolution follows (\ref{eq:frozen}).

Table~\ref{lag:breakers} tests these statements. Each run starts from the same physical fields as the tests of corollary~\ref{cor:singlepoint} and adds one term to GANDALF's right-hand side, integrated with RK4 at $\Delta t=0.00125$; the frozen-line prediction carries the same damping of $g$ where the run has one. The three cases that satisfy (H1) and (H2) approach the prediction as the grid is refined, and the four that violate them leave a difference that does not decrease with resolution. Only this behaviour with resolution tests the theorem, since the sizes of the non-convergent differences depend on the coefficients chosen. In all three convergent cases the decrease from $24^3$ to $32^3$ is small at this time; in the ideal runs of the previous tests the difference in the Hermite spectrum at $t=0.124$ falls again at $48^3$, from $2.4\times10^{-4}$ to $5.2\times10^{-5}$. The runs with viscosity and forcing of $\Phi$ are therefore weak evidence on their own. The proof covers these cases, and so does the test of section~\ref{lag:forcingsec}, in which kicks of $\Phi$ as large as the initial flow leave the mean forced spectrum on the frozen-line prediction to $9.8\times10^{-11}$ on a $48^3$ grid, with spectral convergence in the resolution.

\begin{table}[t]
\centering
\caption{Tests of the hypotheses of theorem~\ref{thm:frozen}. Relative difference $\sum_m|\langle g_m^2\rangle-\langle g_m^2\rangle_{\mathrm{fr}}|/\sum_m\langle g_m^2\rangle_{\mathrm{fr}}$ at $t=0.125$ between an integration of GANDALF's right-hand side with one added term and the frozen-line prediction, on $N^3$ grids ($M=8$, $\Lambda=\sqrt5$, $s=1$). The added terms are a viscosity $\nu_v\nabla_\perp^2\Phi$ in the equation for $\Phi$ with $\nu_v=0.01$; a resistivity $\eta\nabla_\perp^2\Psi$ in the equation for $\Psi$ with $\eta=0.01$; white-noise increments of $\Phi$ or of $\Psi$ with r.m.s.\ gradient $3\,\Delta t^{1/2}$ at every step; and the damping terms $-\eta_gk_\perp^4g$ with $\eta_g=4\times10^{-6}$ and $-\eta_zk_z^2g$ with $\eta_z=0.02$ on every $g_m$.}
\label{lag:breakers}
\begin{tabular}{llccc}
\toprule
added term & hypotheses & $N=16$ & $N=24$ & $N=32$\\
\midrule
none & hold & $2.48\times10^{-3}$ & $2.95\times10^{-4}$ & $2.45\times10^{-4}$\\
viscosity on $\Phi$ & hold & $2.15\times10^{-3}$ & $3.48\times10^{-4}$ & $2.05\times10^{-4}$\\
white forcing of $\Phi$ & hold & $3.10\times10^{-3}$ & $5.44\times10^{-4}$ & $5.02\times10^{-4}$\\
resistivity on $\Psi$ & (H1) fails & $1.88\times10^{-2}$ & $1.88\times10^{-2}$ & $1.89\times10^{-2}$\\
white forcing of $\Psi$ & (H1) fails & $5.09\times10^{-2}$ & $5.07\times10^{-2}$ & $5.06\times10^{-2}$\\
$k_\perp^4$ damping of $g$ & (H2) fails & $3.02\times10^{-2}$ & $3.59\times10^{-2}$ & $3.63\times10^{-2}$\\
$k_z^2$ damping of $g$ & (H2) fails & $3.62\times10^{-2}$ & $4.00\times10^{-2}$ & $4.03\times10^{-2}$\\
\bottomrule
\end{tabular}
\end{table}

At each step GANDALF multiplies $z^\pm$ and every $g_m$ by the same perpendicular and parallel damping factors,
\begin{equation*}
  \exp\big[-\eta\,(k_\perp/k_{\perp,\max})^{2r}\,\Delta t\big]
  \qquad\text{and}\qquad
  \exp\big[-\eta_z\,(k_z/k_{z,\max})^{2r_z}\,\Delta t\big] ,
\end{equation*}
where $r$ and $r_z$ set the order of the damping and $k_{\perp,\max}$ and $k_{z,\max}$ are fixed wavenumbers at the boundary of the 2/3 dealiasing rule. Any GANDALF run with $\eta>0$ or $\eta_z>0$ therefore violates both (H1) and (H2). The routine \texttt{force\_alfven\_modes} adds the same increment to $z^+$ and $z^-$, which forces $\Phi$ only and keeps (H1); \texttt{force\_alfven\_modes\_balanced} forces $z^+$ and $z^-$ independently, which forces $\Psi$ and violates (H1). GANDALF's hypercollisions are diagonal in $m$ and satisfy (H2).

\subsection{Echoes}\label{lag:echoes}

Under (\ref{eq:hyp}), $W(m,t)$ and the field-aligned Hermite flux $\Pi_{m+1/2}(t)$ evolve exactly as in the linear problem $\partial_th=-L_0h$, by (\ref{lag:singlepoint}). Any return of free energy to low $m$ that appears in them must therefore already be present in that linear problem for the given initial data: an initial condition that carries unmixing correlations, which produces a linear echo; the recurrence of a truncated hierarchy; or the slow phase mixing of components with small field-aligned wavenumbers. The random advection by the turbulence does not convert phase mixing into unmixing in $W(m)$. With white forcing the same holds for the mean $\mathrm E\,W(m,t)$, injection by injection (corollary~\ref{cor:forcing}). This is the sense in which ideal, smooth dynamics without spatial dissipation of $g$ cannot produce stochastic echoes in $W(m)$. The statement concerns single-point statistics. It does not say that the flux computed with $\partial_z$, or a split into phase-mixing and unmixing parts by the sign of $k_z$, shows no backward transfer: section~\ref{sec:decorrelation} shows that such diagnostics report backward transfer for pure phase mixing along tilted field lines. Nor does it constrain fluxes resolved in $k_\perp$, or say that the dynamics is linear in the Eulerian frame, in which $W(k_\perp)$ evolves nonlinearly.

Stochastic echoes were studied in drift kinetics with a straight guide field \citep{Schekochihin2016,Parker2016} and in a solvable Vlasov--Poisson model \citep{AdkinsSchekochihin2018}. In both settings the particles stream along a fixed direction, and the nonlinearity can change the wavenumber along it. In drift kinetics the $\boldsymbol E\times\boldsymbol B$ flow, whose stream function $\Phi$ is proportional to the electrostatic potential, depends on $z$, while the streaming is along $\zhat$, so that
\begin{equation*}
  [D_t,\partial_z]\,f=-\{\partial_z\Phi,f\}\neq0 ,
\end{equation*}
and the flow changes the parallel wavenumber of $g$, coupling positive and negative $k_\parallel$. This coupling is the mechanism of the echo. When the advecting flow does not depend on $z$ there is no such coupling and no echo, as noted by \citet[appendix~B]{Schekochihin2016} and \citet[\S5.2]{KanekarThesis2015}. In KRMHD the shear $\partial_z\uE$ that would change $k_\parallel$ also enters the flux equation through $\partial_z\Phi$, and it bends the field lines together with $g$: under (H1) the commutator (\ref{lag:commutator}) vanishes, and by corollary~\ref{cor:kpar} no part of $g$ moves between positive and negative field-aligned wavenumbers. The frozen-in shear (\ref{lag:shear}) is the simplest example: the Eulerian $k_z$ of a mode changes sign while its field-aligned wavenumber stays fixed.

\citet{Meyrand2019} found in simulations of KRMHD turbulence that much of the free-energy flux to high $m$ by phase mixing is cancelled on average by a return flux, and that the Hermite spectrum is a little steeper than $m^{-1}$, against $m^{-1/2}$ for linear phase mixing; they attributed both to stochastic echoes. They described the echo as the scattering of a compressive wave by an Alfv\'en wave propagating in the opposite direction along the field line, which changes the sign of the compressive wave's $k_\parallel$, after which its phase unmixes. \citet{Kawazura2020} measured an $m^{-1}$ spectrum at $\beta_i=0.1$ in the range $k_\perp\rho_i\ll1$ of their gyrokinetic simulations, against $m^{-1/2}$ for linear phase mixing, and read it as an indication that phase mixing is suppressed by stochastic echoes. The theorem does not contradict these measurements. \citet{Meyrand2019} dissipated the Alfv\'enic fields with eighth-order hyperviscosity and hyper-resistivity and the kinetic field with a hyperviscous term $-\mu k_\perp^8g_m$, which violate (H1) and (H2). They also noted that their kinetic equation might seem to be a linear equation in a Lagrangian frame, citing S09, and argued that, because the field develops structure at arbitrarily small scales after about one eddy-turnover time, any resistivity destroys the identities of the field lines and with them this frame. The simulations of \citet{Kawazura2020} also resolve $k_\perp\rho_i\gtrsim1$ and solve a gyrokinetic system, to which the theorem, derived for KRMHD, does not apply directly. The theorem shows that under (\ref{eq:hyp}) the sign of the field-aligned wavenumber cannot change, so an echo of the kind described by \citet{Meyrand2019} requires the non-ideal terms, or a failure of deterministic flux freezing for rough fields in the limit $\eta\to0$.

Such failures of flux freezing have been reported and discussed for magnetohydrodynamic turbulence. \citet{Eyink2013} reported a breakdown of flux freezing in such turbulence at high conductivity. \citet[\S VIII.C]{Eyink2018} argued that in a turbulent inertial range the ideal Ohm's law holds only in a sense that does not imply flux freezing, and on these grounds disputed the flux-freezing statements of S09 and \citet{Kunz2015}. \citet{Jafari2026} argued that magnetic flux is then conserved only statistically, along stochastic path lines. These statements concern rough fields, which (H3) excludes.

Two further measurements of \citet{Meyrand2019} are constrained only in part. Their Hermite flux was resolved in $k_\perp$ shells; it is not a single-point statistic, and the theorem constrains only its sum over shells. They also measured spectra of the compressive fluctuations along traced field lines, found that these spectra track the Alfv\'enic ones, and attributed this parallel cascade to the resistive relabelling of field lines. Corollary~\ref{cor:kpar} is consistent with that attribution: under (\ref{eq:hyp}) no part of $g$ moves between field-aligned wavenumbers. Their runs are also forced, which (\ref{eq:hyp}) excludes, and the theorem does not restrict the field-aligned wavenumbers at which a forcing applied at fixed Eulerian wavenumbers injects $g$.

\citet{Adkins2025} found, in agreement with \citet{Zhou2023}, that echoes play no significant role in the kinetic reduced electron heating model (KREHM) of \citet{ZoccoSchekochihin2011}, and suggested that the passivity of the compressive sector of KRMHD, which does not act back on the Alfv\'enic fields, may be why echoes appear there. The theorem shows that, for smooth solutions without spatial dissipation of $g$, passivity together with ideal induction excludes stochastic echoes in $W(m)$, and in the mean of $W(m)$ when $g$ is forced white in time. Which non-ideal mechanism operates in the KRMHD simulations, and whether its effect survives the limit $\eta\to0$, is left for future work (see also section~\ref{sec:discussion}).
